%% P30 --- Entropia krzyzowa i dywergencja Kullbacka-Leiblera
%%
%% Blizniak: programs/en/P30-cross-entropy-kl.tex. Ramka w ramke, makro w
%% makro, liczba w liczbe; tools/parity.py porownuje oba pliki. Kazda liczba
%% pochodzi z code/p30_cross_entropy_kl.py, ktorego naglowek niesie pelne
%% odczytanie sasiadow.
%%
%% REGULY TLUMACZA (C4, C8, C12 czytaja KOLEJNOSC):
%%   * cyfra zostaje cyfra, a slowo zostaje slowem -- P25 zlapal oba kierunki;
%%   * odwolanie prowadzi zdanie, nigdy nie stoi za swoja matematyka;
%%   * w matematyce \text{} dla przetlumaczonego slowa, nigdy \mathrm{};
%%   * angielski zwrot, ktory po polsku wymaga symbolu, poprawia sie po
%%     STRONIE ANGIELSKIEJ -- P28 zaplacil za to piecioma ramkami;
%%   * terminy ML zostaja angielskie i odmieniaja sie po polsku.

\program[Entropia krzyzowa i KL]{Entropia krzyżowa i dywergencja
  Kullbacka--Leiblera}
\label{prog:P30}
\index{entropia krzyżowa}
\index{dywergencja Kullbacka--Leiblera}

\begin{outcomes}
\outcome{1--6}{powiedzieć, ile kosztuje kodowanie jednego rozkładu kodem
  drugiego, i nazwać dwie części, na które ten koszt się dzieli}
\outcome{7--12}{powiedzieć, dlaczego nadwyżka nigdy nie jest ujemna i kiedy
  wynosi zero}
\outcome{13--18}{policzyć dywergencję w obu kierunkach i powiedzieć, który
  jest który}
\outcome{19--23}{powiedzieć, ile kosztuje brak masy tam, gdzie cel ją ma, i
  ile kosztuje sytuacja odwrotna}
\outcome{24--33}{przewidzieć, które dopasowanie wybierze każdy kierunek, i
  powiedzieć, dlaczego jeden się rozlewa, a drugi zapada}
\outcome{34--38}{powiedzieć, co obiecuje słowo \enquote{odległość} i którą z
  tych obietnic ta wielkość łamie}
\outcome{39--42}{powiedzieć, ile kosztuje symetryczna alternatywa}
\outcome{43--46}{powiedzieć, który kierunek minimalizuje trening na stałym
  zbiorze danych i dlaczego nikt go nie wybrał}
\end{outcomes}

\begin{quiz}
\item Masz kod zbudowany dla jednego rozkładu, a źródło stosuje się do
  innego. Ile kosztuje średnio jeden symbol? \teachesat{1--2}
  \answerto{Tyle, ile wynosi entropia krzyżowa źródła względem rozkładu
  kodu. Jest to własna entropia źródła plus nadwyżka, a o tej nadwyżce jest
  ten program.}
\item Czy cudzy kod może kiedykolwiek pobić ten zbudowany dla twojego
  własnego źródła? \teachesat{7--8}
  \answerto{Nie. Nadwyżka nigdy nie jest ujemna \dash{} to nierówność
  Jensena z Programu~\ref{prog:P19} zastosowana do ilorazu \dash{} i wynosi
  zero dokładnie wtedy, gdy oba rozkłady się zgadzają.}
\item Czy dywergencja od $p$ do $q$ jest taka sama jak dywergencja od $q$ do
  $p$? \teachesat{13--14}
  \answerto{Nie, a na dwóch rozkładach z Programu~\ref{prog:F05} obie
  odpowiedzi różnią się o czynnik $\val{p30.f05.ratio}$.}
\item Ile kosztuje brak masy tam, gdzie cel ją ma?
  \teachesat{19--20}
  \answerto{W jednym kierunku nieskończoność, a w drugim nic, i ta jedna
  asymetria sprawia, że jedna strata się rozlewa, a druga zapada.}
\item Pewną wielkość nazwano odległością. Co ci właśnie obiecano?
  \teachesat{34--35}
  \answerto{Między innymi nierówność trójkąta, którą podaje
  Program~\ref{prog:F09}: żaden objazd nie jest krótszy od drogi wprost.
  Ta wielkość ją łamie, i to nierzadko.}
\end{quiz}

% ============================================================
\section{Ile kosztuje użycie niewłaściwego kodu}
\label{sec:P30-cost}
\index{entropia krzyżowa!jako koszt kodowania}

\begin{fr}
Program~\ref{prog:P29} zostawił cię z kodem i tabelą.

Jego kod zbudowano dla alfabetu $(\frac12, \frac14, \frac18, \frac18)$, gdzie
cztery długości słów kodowych wynoszą $1$, $2$, $3$ i $3$ bity, i daje
średnio dokładnie $\val{p29.code.h}$ bita na symbol \dash{} entropię, bez
żadnego marnotrawstwa.

Teraz wyślij tym samym drutem inne źródło. Symbole przychodzą z
prawdopodobieństwami $(\frac25, \frac15, \frac15, \frac15)$, a kod pozostaje
bez zmian, bo kod jest tym, co ma odbiorca.

Ile kosztuje teraz średnio jeden symbol?
\dotline
\nextframe
\end{fr}

\begin{fr}
\ans{$\val{p30.src.ce}$ bity \dash{} długości ważone nowymi
prawdopodobieństwami}

\[ \frac25 \times 1 + \frac15 \times 2 + \frac15 \times 3
   + \frac15 \times 3 \;=\; \val{p30.src.ce} \]

I to jest całość: \textbf{płacisz długość każdego słowa kodowego tyle razy,
ile razy trafia się jego symbol.} Pisząc $l_i$ na długości, a $p$ na źródło,
ta średnia to $\sum_i p_i l_i$.

A kod zupełny ma $l_i = \logb{2}\frac{1}{q_i}$ dla rozkładu $q$, dla którego
go zbudowano, więc ta sama średnia to

\[ H(p, q) \;=\; -\sum_i p_i \logb{2} q_i \]

co nazywa się \textbf{entropią krzyżową} $p$ względem $q$. To strata z
Programu~\ref{prog:P18}, i to jest właśnie to, co ona mierzy.

A teraz: ile dałby średnio \emph{najlepszy dostępny} kod dla tego drugiego
źródła i gdzie w Programie~\ref{prog:P29} ta liczba już stoi zapisana?
\dotline
\nextframe
\end{fr}

\begin{fr}
\ans{$\val{p30.src.h}$ bita \dash{} własna entropia źródła}

Program~\ref{prog:P29} ją policzył i wydrukował obok długości kodu, bo
entropia jest najkrótszą średnią, jaką może osiągnąć jakikolwiek kod dla tego
źródła.

Płacisz więc $\val{p30.src.ce}$, a najlepsze dostępne to
$\val{p30.src.h}$.

Ile kosztuje cię niewłaściwy kod i gdzie widziałeś \emph{tę} liczbę
wcześniej?
\dotline
\nextframe
\end{fr}

\begin{fr}
\ans{$\val{p30.src.kl}$ bita \dash{} to ta sama luka z
Programu~\ref{prog:P29}}

Tamten program zmierzył dokładnie to i nazwał to \emph{zaokrągleniem długości
do liczby całkowitej}, bo słowo kodowe nie może mieć
$\logb{2}\frac{5}{2}$ bita długości.

Jest tym, i jest jeszcze czymś.

\begin{center}
\emph{Koszt niewłaściwego kodu minus koszt właściwego\\
to wielkość sama w sobie.}
\end{center}

Ma nazwę \dash{} \textbf{dywergencja Kullbacka--Leiblera} rozkładu $p$ od
$q$, zapisywana $\KL{p}{q}$ \dash{} i o niej jest cała reszta tego programu.

\[ H(p, q) \;=\; H(p) \;+\; \KL{p}{q} \]

Sprawdź na liczbach: $\val{p30.src.h} + \val{p30.src.kl}
= \val{p30.src.ce}$.
\end{fr}

\begin{fr}
W tym dojściu jest szczegół, przy którym warto się zatrzymać, bo to on
sprawia, że tożsamość jest tożsamością, a nie przybliżeniem.

Najlepszy kod dla drugiego alfabetu z Programu~\ref{prog:P29} ma długości
$1$, $2$, $3$, $3$, a jego suma Krafta $\sum_i 2^{-l_i}$ wynosi dokładnie
$1$. Liczby $q_i = 2^{-l_i}$ sumują się więc do jedynki, co czyni z nich
rozkład.

Który to rozkład?
\dotline
\nextframe
\end{fr}

\begin{fr}
\ans{$(\frac12, \frac14, \frac18, \frac18)$ \dash{} \emph{drugi} wiersz
tabeli z Programu~\ref{prog:P29}}

Dwa alfabety, które tamten program zestawił, żeby przeciwstawić przypadek
dyadyczny niedyadycznemu, okazują się \textbf{źródłem i kodem jednej
dywergencji}. Nic nowego tu nie policzono. Liczba już była wydrukowana; ten
program dokłada to, czym ona jest.

\begin{note}
To także powód, dla którego kod i rozkład są tym samym obiektem widzianym
dwa razy. Kod zupełny \emph{jest} twierdzeniem o tym, jak prawdopodobny jest
każdy symbol \dash{} daje krótkie słowa symbolom, których się spodziewa
\dash{} a dywergencja mierzy, jak bardzo to twierdzenie było błędne.
\end{note}

\mermaidfig{p30-two-pieces}{Na co dzieli się koszt kodu}{koszt w dwóch
  częściach}

Odejmij dwie sumy, $-\sum_i p_i \logb{2} q_i$ oraz
$-\sum_i p_i \logb{2} p_i$, i zbierz wyrazy. Czym jest dywergencja, wypisana
wprost?
\dotline
\nextframe
\end{fr}

% ============================================================
\section{Nadwyżka nie może być ujemna}
\label{sec:P30-nonneg}
\index{dywergencja Kullbacka--Leiblera!nieujemność}

\begin{fr}
\ans{$\KL{p}{q} = \sum_i p_i \logb{2} \dfrac{p_i}{q_i}$}

Oba logarytmy dzielą wagę $p_i$, więc zbierają się w jeden logarytm ilorazu:
średnia ważona z $\logb{2}$ ilorazu $p_i / q_i$, i nic poza tym.

Oto pytanie, do którego ta definicja zaprasza, i warto na nie odpowiedzieć,
zanim pójdziesz dalej, bo większość ludzi się przy nim waha.

Czy cudzy kod mógłby kiedykolwiek \emph{pobić} ten zbudowany dla twojego
własnego źródła \dash{} czy ta suma mogłaby wyjść ujemna?
\dotline
\nextframe
\end{fr}

\begin{fr}
\begin{ansblock}
Nie, nigdy. A dowód to dwie linijki twierdzenia, które już masz.

Program~\ref{prog:P19} podaje nierówność Jensena: dla funkcji wklęsłej
wartość w średniej jest co najmniej średnią wartości. Logarytm jest wklęsły,
więc

\begin{align*}
  -\KL{p}{q} &\;=\; \sum_i p_i \logb{2}\frac{q_i}{p_i}
              \;\le\; \logb{2} \sum_i p_i \frac{q_i}{p_i} \\
              &\;=\; \logb{2} \sum_i q_i \;=\; \logb{2} 1 \;=\; 0
\end{align*}

a zatem $\KL{p}{q} \ge 0$, zawsze.
\end{ansblock}

Nic nowego nie trzeba było definiować. To nierówność z
Programu~\ref{prog:P19} zastosowana do ilorazu, a ostatni krok mówi tylko
tyle, że rozkład sumuje się do jedynki.
\end{fr}

\begin{fr}
Przypadek równości znaczy tyle samo co nierówność, a Jensen podaje i jego:
funkcja ściśle wklęsła daje równość tylko wtedy, gdy uśredniana wielkość się
nie zmienia.

Tutaj uśrednianą wielkością jest $q_i/p_i$.

Pod jakim dokładnie warunkiem na $p$ i $q$ dywergencja wynosi zero?
\dotline
\nextframe
\end{fr}

\begin{fr}
\ans{Tylko wtedy, gdy to ten sam rozkład}

Iloraz $q_i/p_i$ nie zmienia się dokładnie wtedy, gdy wszystkie są równe
\dash{} a dwa rozkłady w stałym stosunku, z których oba sumują się do
jedynki, są tym samym rozkładem, bo ten stosunek musi wynosić jeden.

\begin{center}
\emph{$\KL{p}{q} = 0$ dokładnie wtedy, gdy $p$ i $q$ się zgadzają,\\
a poza tym jest dodatnia.}
\end{center}

Sprawdzone, a nie zadeklarowane: po wszystkich $\val{p30.simplex.n}$
rozkładach na trzech wynikach o mianowniku $\val{p30.tri.den}$, a więc po
każdej ich uporządkowanej parze, dywergencja nigdy nie jest ujemna i wynosi
zero tylko tam, gdzie oba się zgadzają. Wyczerpująco na tej siatce, co czyni
z tego dowód, a nie świadectwo \dash{} rozróżnienie z
Programu~\ref{prog:P14}.

Ta para faktów to dokładnie to, co obiecano i wykorzystano w
Programie~\ref{prog:P22}.

Cały jego wynik \dash{} że $\beta$ w celu z karą KL jest ceną w nagrodzie na
nat \dash{} potrzebuje tylko tego, że dywergencja jest nieujemną miarą tego,
jak daleko jeden rozkład stoi od drugiego, zerową tylko przy zgodzie.
Powiedziano to w rigour boxie i wskazano ten program.

Ten dług jest teraz spłacony, a wystarczyły do tego definicja i Jensen.

\begin{note}
Zauważ, czego \emph{nie} ustalono. \enquote{Zero tylko przy zgodzie} i
\enquote{nigdy nie ujemna} to dwie z rzeczy, które obiecuje słowo
\emph{odległość}. Nie wszystkie, a §6 jest o tej, która zawodzi.
\end{note}
\end{fr}

\begin{fr}
Jeden wniosek, za darmo, i to ten, który w ogóle czyni z entropii krzyżowej
stratę.

Skoro $H(p, q) = H(p) + \KL{p}{q}$, a dywergencja nigdy nie jest ujemna:

\[ H(p, q) \;\ge\; H(p), \quad\text{z równością dokładnie wtedy, gdy } q = p \]

Jeśli więc utrzymasz $p$ na stałe i będziesz zmieniać $q$, entropia krzyżowa
jest najmniejsza dokładnie wtedy, gdy $q$ jest samym $p$, a jej podłogą jest
własna entropia $p$.

Co to mówi o stracie treningowej, która nie chce zejść poniżej pewnej
liczby, jakkolwiek długo trenujesz?
\dotline
\nextframe
\end{fr}

\begin{fr}
\ans{Podłogą jest własna entropia celu i żaden model nie zejdzie pod nią}

Strata, która się wypłaszcza, nie musi oznaczać modelu, który przestał się
uczyć. Jeśli przewidywana rzecz jest naprawdę niepewna, najlepszy możliwy
model wciąż płaci $H(p)$ na symbol, i każdy bieg przeciwko temu celowi
wypłaszcza się właśnie tam.

\begin{trapbox}
\textbf{\enquote{Strata przestała spadać, więc model przestał się uczyć.}}

Mógł dojść. Entropia krzyżowa jest ograniczona z dołu entropią celu, która
jest własnością \emph{danych}, a nie modelu, a tekst z prawdziwą
wieloznacznością ma podłogę powyżej zera.

Pytanie, które stawia wypłaszczona strata, brzmi: która z dwóch rzeczy
zaszła \dash{} i to dywergencja je rozdziela. Nadwyżka ponad podłogą jest
tym, co zostało do nauczenia, i schodzi do zera przy modelu, który ma rację,
a nie przy modelu, który utknął. Czego nie da się zrobić, to odczytać tego z
samej straty, bo strata to podłoga plus nadwyżka, a podłogi nie ma na
ekranie.
\end{trapbox}
\end{fr}

% ============================================================
\section{Nie jest symetryczna}
\label{sec:P30-asymmetry}
\index{dywergencja Kullbacka--Leiblera!asymetria}

\begin{fr}
Wszystko dotąd byłoby prawdą o odległości. Tu się to kończy.

Program~\ref{prog:F05} zbudował rozkład na czterech tokenach i wydrukował go
w trzech temperaturach. Przy $T = \num{0.5}$ jest szpiczasty \dash{}
$\val{f05.sm.t05.p1}$ na wiodącym tokenie i $\val{f05.sm.t05.p4}$ na
ostatnim. Przy $T = 2$ jest szeroki: $\val{f05.sm.t20.p1}$ i
$\val{f05.sm.t20.p4}$. Napisz $p$ na szpiczasty, a $q$ na szeroki \dash{} to
dwa argumenty z §1 z konkretną parą w środku.

Zwróć uwagę, że oba mają ten sam nośnik: każdy token niesie w obu jakąś
wagę, więc niczego z tego, co nastąpi, nie da się zrzucić na brakujący wynik.
Jakakolwiek asymetria się tu pojawi, nie jest tą, o której mówi §4.

Zanim cokolwiek policzysz, spójrz na te dwie wielkości:

\[ \KL{p}{q} \qquad\text{oraz}\qquad \KL{q}{p} \]

Czy spodziewasz się, że będą sobie bliskie?
\dotline
\nextframe
\end{fr}

\begin{fr}
\ans{Nie są: $\val{p30.f05.fwd}$ przeciwko $\val{p30.f05.rev}$ nata}

Czynnik $\val{p30.f05.ratio}$, na dwóch rozkładach różniących się tylko
ustawieniem temperatury, policzony z dwunastu prawdopodobieństw, które
Program~\ref{prog:F05} zapisał.

\begin{center}
\emph{$\KL{p}{q}$ i $\KL{q}{p}$ to różne liczby,\\
a żadna z nich nie jest odległością między $p$ a $q$.}
\end{center}

Kolejność nie jest konwencją, którą wolno wybrać dla wygody. Nazywa ona to,
o który rozkład pytasz.
\end{fr}

\begin{fr}
Odczytaj oba kierunki z definicji, a powód masz przed sobą.

\[ \KL{p}{q} \;=\; \sum_i p_i \ln\frac{p_i}{q_i} \]

Czynnik $p_i$ poza logarytmem wybiera, \emph{gdzie patrzysz}: wyraz liczy się
o tyle, o ile $p$ mówi, że ten wynik jest ważny. Iloraz w środku wybiera,
\emph{ile płacisz}: jest duży tam, gdzie $p$ jest znacznie większe od $q$.

Który więc rozkład wykonuje którą pracę i który z dwóch jest tym, o którym
należy myśleć jak o \enquote{celu}?
\dotline
\nextframe
\end{fr}

\begin{fr}
\begin{ansblock}
\textbf{Pierwszy} argument waży sumę, więc to jego zdanie o tym, gdzie
patrzeć, jest respektowane. To on jest celem.

\textbf{Drugi} jest tym, który się ocenia, a ocenia się go tylko tam, gdzie
pierwszy położył jakąś wagę.
\end{ansblock}

To cała asymetria w jednym zdaniu, a §4 zamienia ją w cenę.
\end{fr}

\begin{fr}
Teraz liczby na siebie zarabiają. Większym z dwóch kierunków był
$\KL{q}{p}$ \dash{} rozkład \emph{szeroki} na pierwszym miejscu.

Spójrz na ostatnią kolumnę z Programu~\ref{prog:F05}. Przy $T = 2$ czwarty
token niesie $\val{f05.sm.t20.p4}$; przy $T = \num{0.5}$ niesie
$\val{f05.sm.t05.p4}$, co tamten program zmierzył jako wzrost ponad
pięćdziesięciokrotny.

Dlaczego postawienie szerokiego na pierwszym miejscu powiększa odpowiedź?
\dotline
\nextframe
\end{fr}

\begin{fr}
\ans{Bo szeroki kładzie wagę na token, którego szpiczasty niemal się wyrzekł}

Czwarty token wnosi do tego kierunku $\val{f05.sm.t20.p4} \times \ln\frac{
\val{f05.sm.t20.p4}}{\val{f05.sm.t05.p4}}$ \dash{} pokaźną wagę razy
logarytm dużego ilorazu. Odwróć parę, a ten sam token wnosi
$\val{f05.sm.t05.p4}$ razy coś, czyli prawie nic, bo szpiczasty rozkład
ledwie tam zagląda.

\begin{trapbox}
\textbf{\enquote{Na moich danych wyszły mniej więcej tak samo, więc kierunek
nie ma znaczenia.}}

Czasem wychodzą. Dwa alfabety z Programu~\ref{prog:P29} dają
$\val{p30.src.kl}$ w jedną stronę i $\val{p30.src.klrev}$ w drugą \dash{}
dość blisko, by raport cytujący którąkolwiek wyglądał tak samo.

To nie jest uspokojenie, to jest problem. \textbf{Wielkość jednego kierunku
nie mówi nic o wielkości drugiego}, więc para, która przypadkiem zgadza się
na dzisiejszych danych, nie powiedziała ci nic o jutrzejszych. Kierunek
trzeba wybrać na podstawie tego, co znaczy, czyli §4 i §5, a nigdy na
podstawie pomiaru, który wypadł blisko.
\end{trapbox}
\end{fr}

% ============================================================
\section{Ile kosztuje zero, w każdą ze stron}
\label{sec:P30-zero}
\index{dywergencja Kullbacka--Leiblera!a zerowe prawdopodobieństwo}

\begin{fr}
Mechanizm stojący za wszystkim, co dalej, to jeden wyraz tej sumy, i jest to
wyraz, w którym prawdopodobieństwo wynosi zero.

Weź cel $p$, który kładzie wagę na jakimś wyniku, i kandydata $q$, który nie
kładzie tam żadnej: $p_i > 0$ oraz $q_i = 0$.

Ile ten wynik wnosi do $\KL{p}{q}$, a ile wnosi do $\KL{q}{p}$?
\dotline
\nextframe
\end{fr}

\begin{fr}
\begin{ansblock}
Do $\KL{p}{q}$ wnosi $p_i \ln \frac{p_i}{0}$, czyli
\textbf{nieskończoność}.

Do $\KL{q}{p}$ wnosi $0 \times \ln \frac{0}{p_i}$, czyli
\textbf{zupełnie nic} \dash{} tego wyniku po prostu nie ma w sumie.
\end{ansblock}

\begin{center}
\emph{Pominięcie czegoś, co cel ma, jest niewybaczalne w jedną stronę\\
i darmowe w drugą.}
\end{center}

To nie jest różnica zaokrąglenia ani kwestia stopnia. To różnica między
liczbą skończoną a brakiem liczby, i z niej wynika wszystko w §5.

Obie połowy §3 i §4 mieszczą się w pięciu linijkach sesji. Zwróć uwagę, że
zaokrąglenie jest wewnątrz funkcji, a nie nałożone na to, co wypisała, więc
widzisz to, co odda REPL:

\transcript{p30-both-directions}

\mermaidfig{p30-which-is-first}{Co rozstrzyga pierwszy argument}{co jest
  pierwsze}
\end{fr}

\begin{fr}
Implementacja, która napotka ten wyraz, zwykle dołoży małą podłogę do
każdego wpisu, żeby nic nie dzieliło przez zero, a nieskończoność stanie się
dużą liczbą skończoną.

Czy to naprawia zachowanie, czy tylko arytmetykę?
\dotline
\nextframe
\end{fr}

\begin{fr}
\ans{Tylko arytmetykę}

Wyraz tej wielkości nadal dominuje każdy inny wyraz sumy, więc dopasowanie
nadal jest napędzane pokrywaniem nośnika celu. Środek zaradczy ukrywa objaw
i zachowuje zachowanie \dash{} co warto wiedzieć przed §5, bo znaczy to, że
tamtejszy wynik przeżywa każdą jego implementację.

I warto powiedzieć, czym ta nieskończoność \emph{nie} jest. Nie jest
przypadkiem numerycznym ani wadą definicji. To zdanie o kodowaniu z §1
odczytane dosłownie: kod, który nie przypisuje żadnego słowa kodowego
symbolowi, w ogóle nie może tego symbolu wysłać, więc średni koszt wysłania
źródła, które go jednak produkuje, jest nieograniczony. Podłoga z
Programu~\ref{prog:P01} i dyscyplina z Programu~\ref{prog:P02} dotyczą
arytmetyki; to dotyczy samego schematu.
\end{fr}

\begin{fr}
Z tej jednej asymetrii wychodzą dwie nazwy i warto je mieć przed pomiarem, a
nie po nim.

\begin{center}
\begin{tabularx}{\linewidth}{@{}lX@{}}
\toprule
\textbf{W przód}, $\KL{p}{q}$ & cel jest pierwszy, więc luka w jego nośniku
kosztuje $\infty$ \\
\textbf{Wstecz}, $\KL{q}{p}$ & kandydat jest pierwszy, więc luka nie
kosztuje nic, a kosztuje za to masa tam, gdzie cel jej nie ma \\
\bottomrule
\end{tabularx}
\end{center}

Mając tylko tyle i przed jakimkolwiek dopasowaniem: jeśli masz przybliżyć
dwugarbny cel czymś, co ma tylko jeden garb, który kierunek spodziewasz się
zobaczyć rozlany, a który wybierający garb?
\dotline
\nextframe
\end{fr}

% ============================================================
\section{Pokrywanie garbów i szukanie garbu}
\label{sec:P30-modes}
\index{dywergencja Kullbacka--Leiblera!pokrywanie i szukanie garbów}

\begin{fr}
\ans{W przód się rozlewa, wstecz wybiera garb \dash{} a §5 to mierzy}

Rozumowanie to tabela z poprzedniej ramki i nic ponadto. W przód nie stać na
lukę nigdzie tam, gdzie cel ma masę, więc musi pokryć oba garby. Wstecz nie
płaci za lukę nic, a płaci za masę w złym miejscu, więc najtańsze, co może
zrobić, to usiąść na jednym garbie i zignorować drugi.

Te dwa zachowania mają nazwy \dash{} \textbf{pokrywanie garbów} i
\textbf{szukanie garbu} \dash{} i to dla nich istnieje ten program. Reszta
sekcji to pomiar, bo nazwa to nie demonstracja.
\end{fr}

\begin{fr}
Oczywisty sposób, żeby to rozstrzygnąć, to dopasować jednogarbnego kandydata
pod każdym z celów i spojrzeć na dwie odpowiedzi. Program~\ref{prog:P05}
próbował ustalić fakt geometryczny dokładnie w ten sposób i zmierzył
poszukiwanie: procedura utknęła długo przed geometrią, a żadne strojenie by
jej nie uratowało.

Warto powiedzieć dokładnie, co byłoby tu z tym nie tak. Dopasowana odpowiedź
zależy od tego, gdzie ruszyło poszukiwanie, kiedy się zatrzymało i jak
sparametryzowano kandydata \dash{} od trzech rzeczy, które nie mają nic
wspólnego z dwiema porównywanymi dywergencjami. Różnica między dwiema
odpowiedziami mogłaby więc pochodzić z każdej z czterech, a nic na stronie by
nie powiedziało, z której.

Jak więc zmierzyć, którego kandydata woli dany cel, tak żeby odpowiedź nie
była zdaniem o twoim optymalizatorze?
\dotline
\nextframe
\end{fr}

\begin{fr}
\ans{Uczynić zbiór kandydatów skończonym i ocenić każdego z nich}

Nie ma tu żadnego optymalizatora. Kandydaci są skończonym zbiorem i
\textbf{każdy z nich jest oceniany}: $\val{p30.bins}$ przedziałów, trójkątny
kandydat dla każdego środka i każdej połówkowej szerokości, razem
$\val{p30.fam.n}$. Najlepszy jest najlepszym ze wszystkich, co czyni z
wyniku dowód dla tej rodziny, a nie świadectwo o niej \dash{} rozróżnienie z
Programu~\ref{prog:P14}, wykonujące dokładnie tę pracę, przed którą
ostrzega P05.
\end{fr}

\begin{fr}
Cel jest dwugarbny: wysoki garb blisko jednego końca, drugi garb blisko
drugiego i płytka dolina między nimi. Niesie $\val{p30.head.mode2}$ procent
masy na drugim garbie \dash{} to nie szczegół w ogonie, to blisko połowy
rozkładu.

Znaczy jeszcze jedna rzecz i łatwo ją pominąć: \textbf{jest ściśle dodatni w
każdym przedziale.} Nigdzie nie ma dokładnego zera.

Dlaczego to tu znaczy, wobec tego, co właśnie ustaliło §4?
\dotline
\nextframe
\end{fr}

\begin{fr}
\ans{Żeby tylko \emph{jeden} z dwóch kierunków mógł pójść w nieskończoność}

Gdyby cel miał w sobie zera, to kandydat rozlany po nich kładłby masę tam,
gdzie cel jej nie ma \dash{} i kierunek wstecz też byłby nieskończony. Oba
kierunki byłyby zmuszone, asymetria byłaby niewidoczna, a pomiar nie
pokazałby niczego.

\begin{note}
To nie jest przypuszczenie. Taki był pierwszy tutaj wypróbowany projekt i
zawiódł dokładnie tak: przy celu z zerami większość kandydatów była
nieskończona w \emph{obu} kierunkach. Druga próba użyła potem celu, którego
garby były zbyt łagodne, i dywergencja wstecz jednak wolała kandydata
szerokiego, bo wyraz $-H(q)$ w $\sum q \ln \frac{q}{p}$ dominuje, gdy nic
nie jest dość ostre, by warto było na to zapaść.

Obie porażki znaleziono, uruchamiając projekt, a nie rozumując o nim, i
dlatego są tutaj, a nie w przypisie: każda z nich to ten sam mechanizm
widziany z boku.
\end{note}
\end{fr}

\begin{fr}
Zanim padnie odpowiedź. Ilu z $\val{p30.fam.n}$ kandydatów spodziewasz się
mieć \emph{skończoną} dywergencję w przód, a ilu skończoną wstecz?
\dotline
\nextframe
\end{fr}

\begin{fr}
\ans{$\val{p30.fin.fwd}$ w przód i $\val{p30.fin.rev}$ wstecz \dash{}
wszyscy}

Każdy kandydat ma skończoną dywergencję wstecz, bo cel jest wszędzie dodatni
i żaden kandydat nie może położyć masy tam, gdzie cel jej nie ma. Skończoną
w przód ma tylko $\val{p30.fin.fwd}$: reszta jest dość wąska, by zostawić
przedział niepokryty, a jeden niepokryty przedział to $\infty$.

\begin{center}
\emph{Kierunek w przód odrzucił już dwie trzecie kandydatów,\\
zanim któregokolwiek z nich oceniono.}
\end{center}

To nie jest preferencja wyrażona niewielkim marginesem. Większość przestrzeni
poszukiwań jest dla niego po prostu niedostępna, a to, co zostaje, to część
szeroka.
\end{fr}

\begin{fr}
A więc pomiar. Każdy kierunek jest oceniany na każdym z
$\val{p30.fam.n}$ kandydatów i podaje się najlepszego.

\begin{center}
\begin{tabularx}{\linewidth}{@{}lXX@{}}
\toprule
& \textbf{w przód} $\KL{p}{q}$ & \textbf{wstecz} $\KL{q}{p}$ \\
\midrule
wybrany kandydat & najszerszy dostępny & pojedyncza masa punktowa \\
pokrywane przedziały & $\val{p30.fwd.width}$ & $\val{p30.rev.width}$ \\
masa zostawiona na drugim garbie & $\val{p30.fwd.mode2}$ procent & żadna, dokładnie \\
\bottomrule
\end{tabularx}
\end{center}

Cel niesie tam $\val{p30.head.mode2}$ procent. W przód zachowuje część z
tego, a wstecz nie zachowuje \emph{nic} \dash{} i to jest zliczenie, a nie
zaokrąglenie: wybrany kandydat jest zerem w każdym przedziale drugiego
garbu, więc nie ma tu małej liczby wyświetlanej jako brak.
\end{fr}

\begin{fr}
I nie jest to własność tego celu. Tę samą enumerację uruchomiono na
$\val{p30.targets.n}$ różnych dwugarbnych celach \dash{} garby różnej
wysokości, różne odstępy, jedna para prawie równa \dash{} i na każdym z nich
kierunek w przód brał najszerszego dostępnego kandydata, a wstecz brał masę
punktową.

Skrypt to asertuje, a nie raportuje, więc zmiana celu nie może po cichu
zafałszować tej sekcji.

\begin{center}
\emph{Zachowanie należy do tego, który rozkład stoi w liczniku,\\
a nie do kształtu celu.}
\end{center}

Powiedz jednym zdaniem, co cel wstecz zrobił z rozkładem, który miał
modelować, i co jego własna wartość straty o tym mówi.
\dotline
\nextframe
\end{fr}

\begin{fr}
\begin{ansblock}
Wyrzucił garb niosący blisko połowy rozkładu, a jego wartość straty nie mówi
o tym nic: strata jest niska, bo siedzenie na jednym garbie naprawdę jest
najtańszą rzeczą dostępną pod tym celem.

Nie ma tu porażki do wykrycia. Cel został zminimalizowany poprawnie.
\end{ansblock}

\begin{aibox}
Tu dwa kierunki przestają być ciekawostką. Cel wariacyjny i polityka z karą
KL minimalizują oba kierunek wstecz, bo to ten kierunek, którego wartość
oczekiwana jest po próbkach, które da się losować z trenowanej rzeczy. Oba
dziedziczą więc szukanie garbu, a model, który po cichu porzucił garb,
raportuje dla siebie dobrą liczbę.

Destylacja minimalizuje kierunek w przód względem pełnego rozkładu
nauczyciela, co dziedziczy pokrywanie garbów: rozleje ucznia po rzeczach,
które nauczyciel rozważył i odrzucił, zamiast pozwolić mu położyć zero tam,
gdzie nauczyciel go nie położył.

Żadne z tego nie jest błędem i żadnego nie wybrano ze względu na zachowanie
wobec garbów. Przyszło z celem, a tabela z §5 mówi, czym jest.
\end{aibox}

\mermaidfig{p30-two-fits}{Którego kandydata bierze każdy kierunek}{jeden cel,
  dwa dopasowania}
\end{fr}

% ============================================================
\section{To nie jest odległość}
\label{sec:P30-notadistance}
\index{dywergencja Kullbacka--Leiblera!nie jest metryką}

\begin{fr}
Tę wielkość bardzo często nazywa się \enquote{odległością KL}, a §3 złamało
już połowę tego, co to słowo obiecuje: odległość od $p$ do $q$ jest
odległością od $q$ do $p$, a ta nie jest.

Program~\ref{prog:F09} dał ci drugą połowę, dla wektorów, i podał ją jako
twierdzenie z dołączonym warunkiem równości.

Co to było?
\dotline
\nextframe
\end{fr}

\begin{fr}
\ans{Nierówność trójkąta \dash{} żaden objazd nie jest krótszy od drogi
wprost}

$\lVert a - c \rVert \le \lVert a - b \rVert + \lVert b - c \rVert$, z
równością tylko wtedy, gdy objazd leży na drodze wprost.

To nie ozdoba. To własność, która czyni odległość użyteczną w rozumowaniu:
mówi, że ograniczenie dwóch odcinków ogranicza całość, a to właśnie pozwala
w ogóle łańcuchowo składać oszacowania.

Czy dywergencja jej przestrzega?
\dotline
\nextframe
\end{fr}

\begin{fr}
\begin{ansblock}
Nie. Oto trójka, dokładna, na trzech wynikach o mianowniku
$\val{p30.tri.den}$:

\[ \KL{a}{c} = \val{p30.tri.direct}, \qquad
   \KL{a}{b} + \KL{b}{c} = \val{p30.tri.via} \]

Objazd jest \textbf{krótszy} \dash{} o $\val{p30.tri.gap}$ nata, co nie jest
marginalnym naruszeniem, tylko mniej więcej połową wartości wprost.
\end{ansblock}
\end{fr}

\begin{fr}
Łatwo byłoby odłożyć to jako patologiczny narożnik, więc to policzono.

Sprawdzono każdą uporządkowaną trójkę rozkładów na tej siatce:
$\val{p30.tri.total}$ z nich.

Ile spodziewasz się naruszeń nierówności \dash{} garstkę, jakiś procent, czy
większość?
\dotline
\nextframe
\end{fr}

\begin{fr}
\ans{$\val{p30.tri.bad}$ z $\val{p30.tri.total}$ \dash{}
$\val{p30.tri.pct}$ procent}

Blisko jednej czwartej. Naruszanie nierówności trójkąta nie jest dla tej
wielkości przypadkiem brzegowym, tylko \textbf{przypadkiem typowym}, a każde
rozumowanie łańcuchowo składające dwie dywergencje, by ograniczyć trzecią,
nigdy nie miało prawa zadziałać.

\begin{trapbox}
\textbf{\enquote{To wielkość odległościopodobna, więc zwykłe intuicje w
większości się przenoszą.}}

Trzy z nich się nie przenoszą, i to te trzy, z których każdy korzysta. Nie
jest symetryczna (§3), nie przestrzega nierówności trójkąta (tutaj, w
jednej czwartej przypadków) i jest nieograniczona \dash{} jeden niepokryty
wynik posyła ją w nieskończoność, więc żadne przeskalowanie nie zmieści jej
w ustalonym zakresie.

Czym \emph{jest}, to dwa fakty ustalone w §2: nigdy nieujemna i zerowa
dokładnie przy zgodzie. Warto je mieć i nie są odległością. Słowa
\enquote{dywergencja} używa się zamiast \enquote{metryka} właśnie dlatego,
że komuś trzeba to raz powiedzieć.
\end{trapbox}

\mermaidfig{p30-three-promises}{Co obiecuje to słowo i co się utrzymuje}{trzy
  obietnice}
\end{fr}

% ============================================================
\section{Dlaczego po prostu jej nie usymetrycznić}
\label{sec:P30-js}
\index{dywergencja Jensena--Shannona}

\begin{fr}
Oczywistą naprawą jest uśrednienie dwóch kierunków względem punktu
środkowego i ktoś to zrobił. \textbf{Dywergencja Jensena--Shannona} to

\[ \JS{p}{q} \;=\; \tfrac12 \KL{p}{m} + \tfrac12 \KL{q}{m},
   \qquad m = \tfrac12(p + q) \]

Jest symetryczna z konstrukcji. Jest zawsze skończona, bo $m$ jest dodatnie
wszędzie tam, gdzie dodatnie jest $p$ albo $q$. Jej pierwiastek jest nawet
prawdziwą metryką, z nierównością trójkąta włącznie.

Naprawia więc każdą skargę z §3 i §6. Ile to kosztowało?
\dotline
\nextframe
\end{fr}

\begin{fr}
\begin{ansblock}
Jest \textbf{ograniczona}, a to ograniczenie jest kosztem.

Dwa rozkłady, które nie dzielą żadnego nośnika, są od siebie oddalone
dokładnie o $\ln 2 = \val{p30.js.cap}$ nata i o nic więcej, jakkolwiek
daleko od siebie są.
\end{ansblock}

Zmierzone na masach punktowych oddalonych o $1$, $2$, $3$ i
$\val{p30.js.far}$ przedziały: każde z $\val{p30.js.seps}$ oddaleń zwraca
$\val{p30.js.cap}$, co do ostatniej cyfry, którą niesie arytmetyka, podczas
gdy dywergencja jest nieskończona dla wszystkich czterech.

Wygląda to na ciekawostkę o masach punktowych. Kiedy w życiu biegu
treningowego rozkład modelu i cel są na rozłącznych nośnikach?
\dotline
\nextframe
\end{fr}

\begin{fr}
\ans{Na starcie, czyli wtedy, gdy najbardziej potrzeba ci wskazówki, dokąd
iść}

To nasycenie nie jest dziwactwem na skraju. To reżim, w którym bieg
treningowy się zaczyna.

Na początku model rozkłada masę mniej więcej gdziekolwiek, co jest właśnie
przypadkiem rozłącznych nośników. Wielkość symetryczna i ograniczona podaje
tam tę samą liczbę, cokolwiek model zrobi, więc nie niesie \textbf{żadnej
informacji o tym, w którą stronę się ruszyć}.

\begin{center}
\emph{Ograniczoność i użyteczny gradient to ta sama własność\\
czytana z dwoma różnymi znakami.}
\end{center}

Rzecz nieograniczona jest nieograniczona dlatego, że wciąż coś ci mówi w
punkcie, w którym rzecz ograniczona już przestała. Żadna z nich nie jest
lepszą wielkością; są odpowiedziami na różne pytania, a pytaniem jest, czy
potrzebujesz porównać dwie liczby, czy podążać za jedną.
\end{fr}

\begin{fr}
\begin{rigourbox}
Trzy rzeczy są tu podane i nieudowodnione: że $\sqrt{\JS{p}{q}}$ przestrzega
nierówności trójkąta, że dywergencja Jensena--Shannona jest ograniczona z
góry przez $\ln 2$ i że to ograniczenie jest osiągane dokładnie na
rozłącznych nośnikach.

Ostatnia z trzech jest sprawdzona w skrypcie tego programu, przy czterech
oddaleniach, bo to na niej wspiera się rozumowanie. Pierwsze dwie są w
każdym kursie teorii informacji. Nic dalej w tej książce z nich nie
korzysta.
\end{rigourbox}
\end{fr}

% ============================================================
\section{Kierunek, który już minimalizujesz}
\label{sec:P30-empirical}
\index{entropia krzyżowa!a rozkład empiryczny}

\begin{fr}
Zostało jedno pytanie i to ono sprawia, że §5 jest czymś więcej niż uwagą o
dwóch celach, które wybrał ktoś inny.

Trenujesz na stałym zbiorze danych, minimalizując entropię krzyżową. Nic w
tym zdaniu nie wspomina dywergencji, a nikt, kto to wybierał, nie myślał o
pokrywaniu garbów.

Zapisz rozkład empiryczny danych jako $\hat p$ \dash{} proporcje w pliku.
Którą wielkość minimalizuje bieg treningowy?
\dotline
\nextframe
\end{fr}

\begin{fr}
\ans{Dywergencję w przód do rozkładu empirycznego \dash{} tę pokrywającą
garby}

\[ H(\hat p, q) \;=\; H(\hat p) \;+\; \KL{\hat p}{q} \]

czyli tożsamość z §1 z $\hat p$ w środku, a $H(\hat p)$ \textbf{nie zależy
od modelu}. To własność pliku. Minimalizowanie entropii krzyżowej po $q$ i
minimalizowanie dywergencji po $q$ to więc to samo zadanie, różniące się o
stałą, której nikt nie ruszy.

Sprawdzone na trzech modelach względem jednego zbioru danych: tożsamość
zachodzi dokładnie przy każdym z nich, $H(\hat p) = \val{p30.emp.h}$ nata
przez cały czas, a przy modelu zgodnym z danymi nadwyżka wynosi dokładnie
zero. Dla modelu jednostajnego czyta się to tak:

\[ \val{p30.emp.ce.uniform} \;=\; \val{p30.emp.h} \;+\; \val{p30.emp.kl.uniform} \]

gdzie pierwszy składnik po prawej ustala plik, a o drugim decyduje wyłącznie
$q$.
\end{fr}

\begin{fr}
Wybór, który zmierzyło §5, został więc dokonany za ciebie, przez kształt
celu, na długo zanim ktokolwiek go przedyskutował.

\begin{center}
\emph{Każdy bieg, jaki uruchomiłeś na stałym zbiorze danych,\\
minimalizował kierunek pokrywający garby.}
\end{center}

Dlatego model językowy zapytany o coś wieloznacznego raczej asekuruje się
po wszystkich możliwościach, niż wybiera jedną: trenował go cel, który płaci
nieograniczoną cenę za brak masy tam, gdzie dane ją miały, a za rozlewanie
się nie płaci nic.

\begin{note}
To domyka także pętlę otwartą przez Program~\ref{prog:P26}. Tamten program
pokazał, że strata treningowa jest ujemną logarytmiczną wiarygodnością, i
powiedział, że nadwyżka ponad entropią celu \enquote{ma nazwę i należy do
P30}. To dywergencja w przód do $\hat p$, schodzi do zera dokładnie przy
dopasowaniu największej wiarygodności, i dlatego dwa opisy tej samej straty
nigdy się nie kłóciły.
\end{note}
\end{fr}

\begin{fr}
Spójrz wstecz na to, co ten program zrobił z tym, co dostał.

Zdefiniował jedną wielkość, $\sum p \ln \frac{p}{q}$, a wszystko inne było
konsekwencją: jest kosztem minus podłogą, nie może być ujemna, bo logarytm
jest wklęsły, i nie jest symetryczna, bo pierwszy argument wybiera, gdzie
patrzeć.

Jedynym naprawdę nowym faktem jest ten z §4 i jest to pojedynczy wyraz sumy:
zero w złym miejscu kosztuje nieskończoność w jedną stronę i nic w drugą.
Obie nazwy z §5, całe §6 i powód, dla którego bieg treningowy się asekuruje,
wychodzą z tego jednego wyrazu.

Czego \textbf{nie} zrobił, to nie porównał dwóch \emph{zmiennych}. Każda
dywergencja tutaj jest między dwoma rozkładami wypisanymi z góry. W chwili,
gdy zapytasz, ile wiedza o jednej wielkości mówi ci o drugiej \dash{} ile z
etykiety da się przewidzieć z warstwy, powiedzmy \dash{} potrzebujesz
dywergencji między rozkładem łącznym a iloczynem jego brzegowych, a ta ma
własną nazwę i własne pułapki.

Program~\ref{prog:P31} ją definiuje i spędza większość swojej długości na
tym, dlaczego stawiane z jej pomocą tezy zwykle nie są uzasadnione.
\end{fr}

\begin{summarybox}
\sumitem{1--3}{\result{\textbf{Entropia krzyżowa to tyle, ile kosztuje kod}:
  $H(p, q) = -\sum_i p_i \logb{2} q_i$}, czyli długości słów kodowych ważone
  tym, jak często trafia się każdy symbol. To strata z
  Programu~\ref{prog:P18} i to jest to, co ona mierzy.}
\sumitem{4--6}{\result{\textbf{Ten koszt się dzieli}: $H(p, q) = H(p) + \KL{p}{q}$}.
  \enquote{Zaokrąglenie} z Programu~\ref{prog:P29}, wynoszące
  $\val{p30.src.kl}$ bita, \emph{jest} dywergencją, a jego dwa alfabety są
  źródłem i kodem \dash{}
  $\val{p30.src.h} + \val{p30.src.kl} = \val{p30.src.ce}$.}
\sumitem{7--9}{\textbf{\result{Nadwyżka nigdy nie jest ujemna}}, na mocy nierówności
  Jensena z Programu~\ref{prog:P19} dla wklęsłego logarytmu, i wynosi zero
  dokładnie wtedy, gdy oba rozkłady się zgadzają. Sprawdzone po każdej
  uporządkowanej parze z $\val{p30.simplex.n}$ rozkładów na siatce.}
\sumitem{11--12}{\result{Entropia krzyżowa ma więc podłogę, a tą podłogą
  jest entropia \emph{celu}}. Strata, która przestaje spadać, mogła dojść, a nie
  utknąć, i sama strata nie powie ci która.}
\sumitem{13--16}{\textbf{\result{Nie jest symetryczna.}} Na rozkładzie z
  Programu~\ref{prog:F05} oba kierunki to $\val{p30.f05.fwd}$ i
  $\val{p30.f05.rev}$ nata, czynnik $\val{p30.f05.ratio}$. Pierwszy argument
  waży sumę, więc to on jest celem; drugi jest tym, co się ocenia.}
\sumitem{17--18}{A para, która wychodzi prawie równa \dash{}
  $\val{p30.src.kl}$ przeciwko $\val{p30.src.klrev}$ \dash{} nie dowodzi
  niczego: \result{wielkość jednego kierunku nie ogranicza drugiego}.}
\sumitem{19--21}{\result{\textbf{Cały mechanizm to jeden wyraz.} Tam, gdzie cel ma
  masę, a kandydat jej nie ma, w przód płaci się $\infty$, a wstecz nic}. To
  nie różnica zaokrąglenia \dash{} to różnica między liczbą a brakiem
  liczby.}
\sumitem{24--30}{\textbf{\result{W przód pokrywa garby, wstecz szuka garbu}},
  wyliczone po $\val{p30.fam.n}$ kandydatach, bez optymalizatora
  gdziekolwiek. Skończoną wartość w przód ma tylko $\val{p30.fin.fwd}$ z
  nich, przeciwko wszystkim $\val{p30.fin.rev}$ wstecz.}
\sumitem{31--33}{Na celu niosącym $\val{p30.head.mode2}$ procent na drugim
  garbie kierunek w przód zachowuje tam $\val{p30.fwd.mode2}$ procent masy,
  a wstecz nie zachowuje \emph{nic} \dash{} na wszystkich
  $\val{p30.targets.n}$ wypróbowanych celach, więc zachowanie należy do
  kierunku, a nie do celu.}
\sumitem{34--38}{\textbf{\result{To nie jest odległość.}} Nierówność trójkąta, którą
  dał ci Program~\ref{prog:F09}, zawodzi dla $\val{p30.tri.bad}$ z
  $\val{p30.tri.total}$ uporządkowanych trójek \dash{}
  $\val{p30.tri.pct}$ procent, czyli w przypadku typowym \dash{} przy czym
  najgorszy objazd jest krótszy o $\val{p30.tri.gap}$ nata.}
\sumitem{39--42}{\result{\textbf{Usymetrycznienie kosztuje sygnał.} Dywergencja
  Jensena--Shannona jest symetryczna, skończona i ograniczona, a
  ograniczeniem jest $\ln 2 = \val{p30.js.cap}$}: masy punktowe oddalone o
  $1$ i $\val{p30.js.far}$ przedziały są oddalone dokładnie o tyle, a to
  reżim, w którym bieg się zaczyna.}
\sumitem{43--45}{\result{\textbf{Trening na stałym zbiorze danych minimalizuje
  dywergencję w przód} do rozkładu empirycznego}, bo
  $H(\hat p) = \val{p30.emp.h}$ nata nie zależy od modelu. Nikt nie wybrał
  kierunku pokrywającego garby; przyszedł on z celem.}
\end{summarybox}

\canyou

\begin{testexercises}
\item Źródło stosuje się do $(\frac12, \frac12)$, a kodujesz je schematem
  zbudowanym dla $(\frac34, \frac14)$. Podaj entropię krzyżową, entropię i
  dywergencję, w bitach.
  \answerto{Długości kodu to $\logb{2}\frac43 = \num{0.415}$ oraz
  $\logb{2}4 = 2$, więc entropia krzyżowa wynosi
  $\frac12(\num{0.415}) + \frac12(2) = \num{1.207}$ bita. Entropia to
  $1$ bit. Dywergencja to różnica, $\num{0.207}$ bita \dash{} tyle kosztuje
  cię niewłaściwy kod na symbol.}
\item Nie licząc niczego, powiedz, która z $\KL{p}{q}$ i $\KL{q}{p}$ jest
  nieskończona, gdy $q$ jest zerem na wyniku, któremu $p$ nadaje wagę, i
  dlaczego.
  \answerto{$\KL{p}{q}$. Pierwszy argument waży sumę, więc ten wynik jest w
  niej z wagą $p_i > 0$, a iloraz $p_i / 0$ jest nieograniczony. W drugim
  kierunku wynik wchodzi z wagą $q_i = 0$ i wypada z sumy zupełnie.}
\item Kolega raportuje, że na jego danych oba kierunki zgadzają się do
  trzech cyfr, i wnioskuje, że dla jego problemu wybór nie ma znaczenia. Co
  jest złego w tym wniosku?
  \answerto{Zgodność jest faktem o tej parze rozkładów, a nie o samej
  wielkości. Dwa alfabety z Programu~\ref{prog:P29} też zgadzają się blisko
  \dash{} $\val{p30.src.kl}$ przeciwko $\val{p30.src.klrev}$ \dash{}
  podczas gdy cele z §5 różnią się wszystkim. Gorzej, wybór rządzi
  \emph{zachowaniem} podczas dopasowywania, więc znaczy najbardziej tam,
  gdzie tych dwóch liczb jeszcze nie policzono.}
\item Destylujesz nauczyciela do mniejszego ucznia i minimalizujesz
  dywergencję z nauczycielem na pierwszym miejscu. Czego spodziewasz się po
  uczniu wobec opcji, której nauczyciel nadał małe, ale niezerowe
  prawdopodobieństwo?
  \answerto{Że zostawi na niej trochę masy. Nauczyciel na pierwszym miejscu
  to kierunek w przód, więc uczeń płaci nieograniczoną cenę za położenie
  zera tam, gdzie nauczyciel go nie położył \dash{} czyli pokrywa garby:
  uczeń rozlewa się po rzeczach, które nauczyciel rozważył i odrzucił,
  zamiast się zdecydować.}
\item Podaj jednozdaniowy powód, dla którego ograniczona, symetryczna
  dywergencja jest kiepskim celem treningowym, mimo że naprawia każdą
  skargę z §3 i §6.
  \answerto{Bo się nasyca: na rozłącznym nośniku, od którego bieg startuje,
  podaje $\ln 2$ cokolwiek model zrobi, więc nie mówi nic o tym, w którą
  stronę się ruszyć.}
\item Strata treningowa wypłaszczyła się na liczbie wyraźnie powyżej zera, a
  zespół spiera się, czy model przestał się uczyć. Która wielkość to
  rozstrzyga i czy da się ją odczytać ze straty?
  \answerto{Dywergencja do celu \dash{} entropia krzyżowa minus własna
  entropia celu. Ze straty jej nie odczytasz, bo strata to podłoga plus
  nadwyżka, a podłoga jest własnością danych, której nie ma na ekranie.
  Trzeba ją oszacować osobno albo ograniczyć.}
\end{testexercises}

\begin{furtherproblems}
\item Pokaż, że $H(p, q) \ge H(p)$ dla każdego $q$, z równością tylko przy
  $q = p$, korzystając wyłącznie z wyniku z §2.
  \answerto{$H(p, q) = H(p) + \KL{p}{q}$, a dywergencja jest nieujemna, więc
  $H(p, q) \ge H(p)$; równość wymaga $\KL{p}{q} = 0$, co według §2 zachodzi
  tylko wtedy, gdy $q = p$.}
\item Kandydat $q$ jest jednostajny na $n$ wynikach. Pokaż, że $\KL{p}{q}$
  wynosi $\logb{2} n - H(p)$ w bitach, i powiedz, co to znaczy o tym, który
  $p$ jest najdalej od jednostajnego.
  \answerto{Każde $q_i$ wynosi $\frac1n$, więc
  $\sum_i p_i \logb{2}\frac{p_i}{1/n} = \sum_i p_i \logb{2} p_i
  + \logb{2} n = \logb{2} n - H(p)$. Dywergencja od jednostajnego to
  dokładnie niedobór entropii, więc najdalej od jednostajnego jest rozkład o
  najmniejszej entropii \dash{} masa punktowa, przy $\logb{2} n$ bitach.}
\item Weź $p = (\num{0.5}, \num{0.5})$ oraz $q = (1, 0)$. Policz oba
  kierunki i obie wartości Jensena--Shannona.
  \answerto{$\KL{p}{q}$ jest nieskończona: $q$ jest zerem tam, gdzie $p$ ma
  wagę. $\KL{q}{p} = 1 \times \ln\frac{1}{\num{0.5}} = \ln 2$. Dywergencja
  Jensena--Shannona wynosi $\val{p30.js.overlap}$ w obu kierunkach, będąc
  symetryczną \dash{} i jest wyraźnie \emph{poniżej} ograniczenia, bo te dwa
  rozkłady jednak dzielą nośnik: $q$ kładzie całą swoją wagę na tym wyniku, na
  którym $p$ ma jej połowę. Sekcja 7 osiąga $\val{p30.js.cap}$ na masach
  punktowych, które nie dzielą niczego, i to jest warunek, a nie poboczny
  szczegół.}
\item Nierówność trójkąta zawodzi dla $\val{p30.tri.pct}$ procent
  uporządkowanych trójek na siatce z §6. Wyjaśnij, dlaczego liczenie trójek
  \emph{uporządkowanych} jest tym, co należy policzyć.
  \answerto{Bo wielkość nie jest symetryczna, więc $(a, b, c)$ i
  $(c, b, a)$ to różne twierdzenia i jedno może zawieść, gdy drugie się
  utrzymuje. Liczenie trójek nieuporządkowanych po cichu zlałoby oba i
  podało niższy odsetek naruszeń dla wielkości, której asymetria jest tu
  sednem.}
\item Załóżmy, że przed policzeniem dywergencji w przód dodajesz małą
  podłogę $\varepsilon$ do każdego wpisu kandydata, tak by żaden wyraz nie
  był nigdy nieskończony. Powiedz, co się zmienia, a co nie.
  \answerto{Zmienia się arytmetyka: niepokryty wynik kosztuje teraz
  $p_i \ln \frac{p_i}{\varepsilon}$, czyli skończenie. Zachowanie się nie
  zmienia: ten wyraz nadal rośnie bez ograniczenia, gdy $\varepsilon$
  maleje, i nadal dominuje sumę, więc dopasowanie nadal jest pchane do
  pokrycia nośnika celu. Podłoga usuwa objaw i zachowuje pokrywanie garbów.}
\end{furtherproblems}
